\documentclass[10pt,a4paper]{amsart}
\usepackage[margin=19mm]{geometry}
\usepackage{amsmath,amssymb,mathtools}
\usepackage[scr=rsfs,cal=euler]{mathalfa}
\usepackage{xfrac}
\usepackage[unicode,hidelinks,bookmarks=false]{hyperref}
\newcommand{\ie}{i.e.\ }
\newcommand{\cf}{cf.\ }
\newcommand{\resp}{respectively\ }
\newcommand{\Subsection}{\subsection}
\providecommand{\nobreakdash}{\nobreak\hbox{-}}
\setlength{\emergencystretch}{2em}
\multlinegap0pt
\usepackage{amssymb}
\usepackage{bm}
\usepackage{csquotes}
\usepackage{xcolor}
\usepackage{mathtools}
\newtheorem{theorem}{Theorem}
\newtheorem{corollary}{Corollary}
\newtheorem{claim}{Claim}
\usepackage{cleveref}
\crefname{theorem}{Theorem}{Theorems}
\crefname{corollary}{Corollary}{Corollarys}
\crefname{claim}{Claim}{Claims}
\title{New Upper bounds on the Mondrian Art Problem}
\author{Thomas Garrison, Chris Seiler, Aliaksei Semchankau}
\date{Source: arXiv:2609.01998v1 (CC BY 4.0)}
\begin{document}
\maketitle

\noindent\emph{Source and licence.} New Upper bounds on the Mondrian Art Problem. Version arXiv:2609.01998v1 (CC BY 4.0), distributed by arXiv under the Creative Commons Attribution 4.0 International licence, http://creativecommons.org/licenses/by/4.0/, as recorded in the arXiv licence metadata for this exact version and shown on the abstract page https://arxiv.org/abs/2609.01998v1. Full licence notice: \texttt{licenses/arxiv-2609.01998-CC-BY-4.0.txt}.\par\smallskip

\noindent\textit{Editorial scope.} Counted unit: the article's theorem thm:main (source0.tex lines 86--90). The article's complete original proof of it is delivered with it (source0.tex lines 172--214, one \texttt{proof} environment of 43 lines, which the article places away from the statement). No mathematics is written, repaired or re-derived: the statement and the proof are the article's own lines, in the article's own notation, and no retained line is substituted or re-flowed.  Retained uncounted as the source-local context the proof needs: lines 99--101.  Nothing was omitted from any retained span, so the package carries no omission record.  No retained line contains a citation, so the wrapper carries no bibliography and no bibliography entry.  No retained line contains a figure environment, an image-inclusion command or drawing code, so no image is delivered and the package declares no asset dependency.  Editorial formatting only, in addition: an AMS \texttt{amsart} preamble that restates the article's own declarations and macros used by the retained lines, namely the environments \texttt{theorem}, \texttt{corollary}, \texttt{claim} on separate counters, together with the standard packages those lines need.  The retained environments print in this document as Theorem 1, Corollary 1, Claim 1 in order of appearance, whereas the article prints the same items with the numbering of its own full document.  The complete native source of the article as distributed in the arXiv e-print for this exact version is bundled unchanged as \texttt{sources/209158/source0.tex}.
\begin{corollary}\label{cor:range}
 For any number $N$ in the range $m(\sqrt{2m}+1) \leq N \leq 2m\ell - m(\sqrt{2m}+1)$ there is a subset $R_N \subseteq \{m - \ell, \ldots, m + \ell\}$ which sums up to $N$.    
\end{corollary}

\begin{theorem}\label{thm:main}

There exists $n_0$ such that for $n\geq n_0$, there exists a Mondrian partition of the $n\times n$ square with defect  $d\leq 10 n^\frac{5}{6}$.

\end{theorem}

\begin{proof}[Proof of \cref{thm:main}]
In this proof we use four constants: $C$, the constant in front of the defect; $c$, the constant for the tile whose area we are approximating; $c_1, c_2$, the constants representing the smallest and largest strip widths.
Let $d := Cn^{5/6}$ be half of the max allowable defect ($C$ is permitted to be a large constant).

Let $T := cn^{7/6}$ be an area we will try to approximate ($c$ is a small constant). All areas will be in the $[T - d, T + d]$ range giving a defect of $2d$.

Now we partition the rectangle $n \times n$ into sub-rectangles $w_1 \times n$, $w_2 \times n$, $\ldots$ such that all $w_1, w_2, \ldots$ are distinct and come from the range $[c_1 n^{1/2}, \ldots, c_2 n^{1/2}]$. To make this partition cover the whole range, we just need 
\[
\sum_{c_1\sqrt{n}}^{c_2\sqrt{n}}w\geq n\ 
\Leftrightarrow\ 
\frac{c_2^2n}{2}-\frac{c_1^2n}{2}\geq n+o(n)\ 
\Leftrightarrow\ 
c_2^2 - c_1^2 \geq 2 + o(1).
\]

Each of the strips $w_i \times n$ we will partition into rectangles $w_i \times h_{i1}, w_i \times h_{i2}, \ldots$ in such a manner that  $w_i < h_{i1} < h_{i2} < \ldots$ holds, so as to ensure that all dimensions across different strips are distinct. Explicitly, take any two rectangles with dimensions $(w_i,h_{ij})$ and $(w_{i'},h_{i'j'})$; even if they have the same widths, their heights must be different so the rectangles have different dimensions. We also have $w_i \ll h_{ij}$ since $w_i=\Theta(n^{1/2})$ and $h_{ij}=\Theta(n^{2/3})$, so the heights and widths can never be equal, so we cannot have rectangles with height and width dimensions swapped. We also want all $w \times h_{ij}$ to have areas in $[T - d, T + d]$.

Given a strip $w_i \times n$, let us write $w = \alpha n^{1/2}$, $c_1 \leq \alpha \leq c_2$. 
Then we want all sides $h_{ij}$ to be in the interval $[T/w_i \pm d/w_i]$. In terms of the lemma above, we have $m = T / w_i \sim \frac{c}{\alpha}n^{2/3}, \ell = d / w_i \sim \frac{C}{\alpha}n^{1/3} $.

Notice, that as $T/w_i - d/w_i \gg n^{2/3} \gg n^{1/2}$, inequalities $w_i < h_{i1} < \ldots$ hold automatically.
\begin{claim}\label{clm:constants}
     We can choose constants $c,c_1,c_2,C$ such that all the numbers from $m\sqrt{2m}\leq n \leq 2m\ell - m\sqrt{2m}$ can be covered by the steps in our algorithm.   
\end{claim}

By \cref{cor:range}, we can pick heights for strips of width $w_i$ that sum to $n$ if $n$ lies in the range $[m(\sqrt{2m}+1), 2m\ell - m(\sqrt{2m}+1)]$. Since the $+1$ adds a term that is at most $O(n^{2/3})$, it is enough to show $n$ lies in $[m\sqrt{2m}, 2m\ell - m\sqrt{2m}]$ up to $o(n)$. Substituting in for $m =  \frac{c}{\alpha}n^{2/3}, \ell =  \frac{C}{\alpha}n^{1/3} $ this range becomes $ [\sqrt{2}c^{3/2}\alpha^{-3/2}n, (2Cc\alpha^{-2} - \sqrt{2}c^{3/2}\alpha^{-3/2})n]$. To have number $n = 1 \cdot n$ contained in this range we need inequalities
\[
\sqrt{2}\left(\frac{c}{c_1}\right)^{3/2} 
\leq 
1 
\leq 
2\frac{Cc}{c_2^2} - \sqrt{2}\left(\frac{c}{c_2}\right)^{3/2}
\]
and
\[
c_2^2 - c_1^2 \geq 2 + o(1)
\]
to hold.

We can do that. For instance, our inequalities are satisfied with constants $c = 1, c_1=2, c_2=2.5 ,  C = 5$.

This completes the proof.
\end{proof}

\end{document}
